\documentclass{article}
\usepackage{hyperref}
\title{A health callback received during a switch is not yet consumed}
\begin{document}
\maketitle
The pinned ARIAC switcher can receive a second fault report while
waiting on an earlier activation request, retain the report's values,
and lose the bit that tells its main loop to act. This happens even
if the fault injector correctly publishes occurrence-scoped health
and the old switch completes exactly as requested. Repairing the
fault injector alone cannot repair the consumer.

\paragraph{Executed cut.}
The native code starts both groups at tick 1. The service stub supplies
a floor-only disable callback at tick 2 during the wait, and returns the
old switch's success at tick 4. The callback sets the pending bit; the
old switch handler clears it at the end of its iteration. After the
wait, a nominal spin that does not deliver a second health message
leaves the floor enabled in the mock plant. If an identical health
message arrives at tick 9, the unmodified switcher finally requests
the floor disable then. In the control where that report is first
delivered after completion, the disable is requested at receipt.
In another control the report clears inside the old wait, before the
next decision, so a new disable is no longer warranted by current health.

The selected call orders are exercised against the exact Python
switcher blob with subscription and service stubs. They do not
establish that a live executor chose these orders or that these
ticks are physical times. In the checked ROS Humble implementation,
\texttt{spin\_until\_future\_complete} executes ready callbacks while
waiting. Its source permits the illustrated cut, but does not imply
all reports arrive before the future. ARIAC's nominal 10 Hz simulation
publication setting cannot by itself guarantee another delivered
callback by a service deadline.

\paragraph{Scope of the adjustment.}
With one thread, consume the pending bit just before taking the local
snapshot that will drive the request, and do not block on a new spin
when a newer report is already pending. That partitions received
callbacks into those covered by the old snapshot and those still
pending; late completion of the old request cannot acknowledge the
new ones. If there were multiple possible concurrent writers this
would require atomic snapshot-and-consume or a checked sequence
comparison. In the matched diagnostic, the adjusted switcher
requests the floor disable at tick 4 in both persistent cases and
sends no stale disable when the input cleared during the wait. The
localized adjustment agrees with a direct latest-report reference
in 341 replayed callback words of length at most four as well as
the four named examples. Under successful exact group switches and
this snapshots-only control law, the same argument applies to
any finite during-wait sequence: the last snapshot at the cut
fully specifies the requested next group state. The replay does
not make the trigger path faster than the old wait.

A failed old service still leaves the unmodified group state
optimistically set in the adjusted variant. A repeated healthy report
then still does not retry that request. Independent evidence of
the actual controller state remains necessary in that case.
And when protection must take effect before the old wait finishes,
the reported latest value is not a substitute for an effect-side
hold. Those failures call for an authority or execution change,
not a mailbox adjustment. This is an I1 source calibration,
not a comparison win against a correctly versioned executive or
evidence of completed physical recovery.
\end{document}
